\subsection{Module Homomorphisms, Kernels, and Images}

Throughout this subsection, $R$ is a nonzero commutative ring with identity $1_R$.

\begin{definition} \label{def:module-homomorphism} \hlt{Module Homomorphism}\\
Let $M,N$ be $R$-modules and let $\tau:M\to N$.
\begin{enumerate}[label=(\roman*)]
\setlength{\itemsep}{0pt}
\item The map $\tau$ is \hlt{$R$-linear} if $(\forall\,(r,s,u,v)\in R^2\times M^2)\,\tau(ru+sv)=r\tau(u)+s\tau(v)$. An $R$-linear map is an \hlt{$R$-homomorphism}.
\item The set of all $R$-homomorphisms $M\to N$ is $\operatorname{Hom}_R(M,N)$.
\item An $R$-homomorphism $M\to M$ is an \hlt{endomorphism}.
\item If $\tau$ is an $R$-homomorphism, then:
\begin{enumerate}[label=(\alph*)]
\setlength{\itemsep}{0pt}
\item $\tau$ is a \hlt{monomorphism} $\Longleftrightarrow\tau$ is injective;
\item $\tau$ is an \hlt{epimorphism} $\Longleftrightarrow\tau$ is surjective;
\item $\tau$ is an \hlt{isomorphism} $\Longleftrightarrow\tau$ is bijective.
\end{enumerate}
\end{enumerate}
\end{definition}

\begin{proposition} \label{prop:hom-module-kernel-image} \hlt{Hom, Kernel, and Image}\\
Let $M,N$ be $R$-modules with additive identities $0_M,0_N$, and let $\tau\in\operatorname{Hom}_R(M,N)$.
\begin{enumerate}[label=(\roman*)]
\setlength{\itemsep}{0pt}
\item For $(\sigma,\rho,m)\in\operatorname{Hom}_R(M,N)^2\times M$, set $(\sigma+\rho)(m)\defeq\sigma(m)+\rho(m)$.\\
For $(a,\sigma,m)\in R\times\operatorname{Hom}_R(M,N)\times M$, set $(a\sigma)(m)\defeq a\sigma(m)$. These pointwise operations make $\operatorname{Hom}_R(M,N)$ an $R$-module.
\item $\ker\tau\defeq\tau^{-1}(\{0_N\})\leq M$.
\item $\operatorname{im}\tau\defeq\tau(M)\leq N$.
\item $\tau$ is a monomorphism $\Longleftrightarrow\ker\tau=\{0_M\}$.
\end{enumerate}
\begin{proof}
$(\forall\,a\in R)\,a0_N=0_N$ by additive cancellation in $a0_N=a(0_N+0_N)=a0_N+a0_N$.\\
$\tau(0_M)=0_N$ by Definition \ref{def:module-homomorphism}(i) with $(r,s,u,v)=(1_R,1_R,0_M,0_M)$ and additive cancellation in $N$.
\begin{enumerate}[label=(\roman*)]
\setlength{\itemsep}{0pt}
\item[$(i)$] Set $\mathcal{F}(M,N)\defeq\{f:f:M\to N\}$ and equip it with the pointwise operations in the statement.\\
Zero map $0_{\mathcal{F}}:M\to N$, $m\mapsto0_N$, is additive identity. Module axioms hold pointwise in $N$ and in $\mathcal{F}(M,N)$ by Assumption \ref{assump:function-extensionality}.\\
Thus $\mathcal{F}(M,N)$ is an $R$-module by Definition \ref{def:module-over-ring}. The zero map is $R$-linear by $(\forall\,r,s\in R)\,r0_N+s0_N=0_N$.\\
Fix $\sigma,\rho\in\operatorname{Hom}_R(M,N)$, $a,b,r,s\in R$, and $u,v\in M$.\\
$(a\sigma+b\rho)(ru+sv)=r(a\sigma+b\rho)(u)+s(a\sigma+b\rho)(v)$ by $R$-linearity of $\sigma,\rho$ and ring commutativity.\\
Hence $a\sigma+b\rho\in\operatorname{Hom}_R(M,N)$ by Definition \ref{def:module-homomorphism}(i). Therefore $\operatorname{Hom}_R(M,N)\leq\mathcal{F}(M,N)$ by Theorem \ref{thm:submodule-test}.
\item[$(ii)$] $0_M\in\ker\tau$ by $\tau(0_M)=0_N$. For $(a,b,x,y)\in R^2\times(\ker\tau)^2$, $\tau(ax+by)=a\tau(x)+b\tau(y)=a0_N+b0_N=0_N$.\\
Thus $\ker\tau\leq M$ by Theorem \ref{thm:submodule-test}.
\item[$(iii)$] $0_N\in\operatorname{im}\tau$ by $0_N=\tau(0_M)$. Fix $(a,b,y_1,y_2)\in R^2\times(\operatorname{im}\tau)^2$.\\
Choose $x_1,x_2\in M$ with $(y_1,y_2)=(\tau(x_1),\tau(x_2))$ by the definition of the image.\\
Then $ay_1+by_2=\tau(ax_1+bx_2)\in\operatorname{im}\tau$ by Definition \ref{def:module-homomorphism}(i). Thus $\operatorname{im}\tau\leq N$ by Theorem \ref{thm:submodule-test}.
\item[$(iv)$] ($\Rightarrow$) Assume $\tau$ is a monomorphism, hence injective by Definition \ref{def:module-homomorphism}(iv).\\
$(\forall\,m\in\ker\tau)\,m=0_M$ by injectivity and $\tau(m)=0_N=\tau(0_M)$. Therefore $\ker\tau=\{0_M\}$ by $\ker\tau\subseteq\{0_M\}$ and (ii).
\item[$(iv)$] ($\Leftarrow$) Assume $\ker\tau=\{0_M\}$ and fix $m,m'\in M$ with $\tau(m)=\tau(m')$. Then $\tau(m-m')=0_N$ by Definition \ref{def:module-homomorphism}(i).\\
Thus $m-m'\in\ker\tau=\{0_M\}$ and $m=m'$. Thus $\tau$ is injective and hence a monomorphism by Definition \ref{def:module-homomorphism}(iv).
\end{enumerate}
\hlt{Structural Commentary}: Hom closes pointwise; kernels use differences; images use lifted representatives.\\
\hlt{Key Insight}: $R$-linearity transports combinations and detects equality through kernel differences.
\end{proof}
\end{proposition}
